\subsection{Resolving an open problem in a special case}
\label{sec:new-mathematics}
\emph{Mistake-Bounded Language Generation} \citep[][P29]{mistake}
studies generators that learn from positive examples and output an unobserved
point at each step. A mistake occurs when the output is outside the target
language. The paper asks how the total number of mistakes trades off against
how quickly the generator stops making them (Section~8, Open Direction~(2)).
We resolve this question for deterministic generation on the family underlying
the paper's tradeoff construction. GenLimitLib's separate definitions of
mistake and convergence guarantees helped guide this analysis.
As in P29, observations are distinct, and time counts the points observed so far.

We consider the family of languages $L_i=P_i\cup R_i$,
where $P_i=\bigcup_{k\le i}C_k$. The blocks $C_i$ are nonempty and finite,
the sets $R_i$ are countably infinite, and all these sets are mutually disjoint.
Thus each successive language contains one more finite block, together with
an infinite part $R_i$ that belongs only to that language. We call this the \emph{staircase family} and write $s_i=|P_i|$.
We measure convergence by the number of distinct observations needed before
all later outputs are guaranteed to be valid, for every presentation order.

For this family, the tradeoff comes from a simple choice. When the observed
set is exactly $P_k$, the generator can output in $R_k$, which is valid only
for $L_k$, or in $C_{k+1}$, which is valid for every later target.
Away from these points, it can choose an unobserved point common to all
compatible targets. Let $\sigma_k=1$ record the choice of $R_k$ and
$\sigma_k=0$ the choice of $C_{k+1}$. Write $M_i(G)$ for the worst-case
number of mistakes of a generator $G$ on $L_i$, and $D_i(G)$ for its
worst-case convergence time.

\begin{result}[Optimal mistake and convergence guarantees]
\label{thm:staircase-frontier}
For a binary sequence $\sigma=(\sigma_k)_{k\ge1}$, let
$J_i(\sigma)=\{1\le k<i:\sigma_k=1\}\cup\{i:\sigma_i=0\}$ and define
\[
 m_i(\sigma)=|J_i(\sigma)|,\qquad
 d_i(\sigma)=\max\bigl(\{0\}\cup\{s_k+1:k\in J_i(\sigma)\}\bigr).
\]
Each sequence $\sigma$ is realized by a deterministic generator $G_\sigma$
with $M_i(G_\sigma)=m_i(\sigma)$ and $D_i(G_\sigma)=d_i(\sigma)$ for all $i$.
Conversely, every deterministic generator $G$ admits a sequence $\sigma$
such that $m_i(\sigma)\le M_i(G)$ and $d_i(\sigma)\le D_i(G)$ for every $i$.
These are exactly the Pareto-optimal guarantees: none can be improved
without worsening another.
\end{result}

The theorem gives matching constructions and lower bounds for every optimal
tradeoff on this family. It determines both the number and the latest possible
timing of mistakes, for arbitrary positive finite block sizes
(Appendix~\ref{app:p29-frontier}).
